%% kapa: exam=1 round=20 year=2009 period=1 form=A subject=civil_law subject_ko=민법 seq=14 number=14
%% source: Q-Net 감정평가사 기출문제 (cst00309) artlSeq=5024178 file=extracted/1008/5024178/2009년 제20회 감정평가사 1차시험 기출문제/1교시.hwp
%% conversion: kordoc 4.24.51 -> tools/gen1.py (auto); review=proofread-sonnet; flags=none
\begin{question}{14}
등기의 추정력에 관한 설명으로 옳지 \uline{않은} 것은? (다툼이 있으면 판례에 의함)
\begin{choices}
  \item 소유권이전등기는 그 추정력을 다투는 측에서 무효사유를 주장ㆍ입증하지 않는 한, 등기원인 사실에 관한 입증이 부족하다는 이유로 그 등기가 무효라고 단정할 수 없다.
  \item 소유권이전청구권 보전을 위한 가등기가 경료되어 있으면 소유권이전등기를 청구할 어떤 법률관계가 있다고 추정된다.
  \item 전(前) 소유자가 사망한 이후에 그 명의의 신청에 의하여 이루어진 이전등기는 일단 원인무효의 등기라고 볼 것이어서, 특별한 사정이 없는 한 등기의 추정력이 인정되지 않는다.
  \item 소유권이전등기가 경료되어 있는 경우에 그 등기명의자는 제3자에 대해서 뿐만 아니라 그 전(前)소유자에 대해서도 적법한 등기원인에 의하여 소유권을 취득한 것으로 추정된다.
  \item 소유권이전등기가 원인 없이 말소된 경우에는 그 회복등기가 경료되기 전이라 하더라도 말소된 등기의 최종명의인은 적법한 권리자로 추정된다.
\end{choices}
\end{question}
\answer{2}
